\documentclass[10pt,letterpaper]{article}
\usepackage{cornell-notes}

\title{Clause 33.05.08.02: Operations On Atomic Types}
\author{}
\date{}
\setDocTitle{Clause 33.05.08.02: Operations On Atomic Types}
\setDocSubtitle{C++ 2024}
\setDocOwner{C++ 2024 Cornell Notes}
\setDocDate{\today}
\setCornellCollection{C++ 2024 Cornell Notes}
\setCornellUnitType{Clause}
\setCornellUnitNumber{33.05.08.02}
\setCornellUnitTitle{Operations On Atomic Types}
\hypersetup{pdftitle={Clause 33.05.08.02: Operations On Atomic Types},pdfsubject={Cornell notes},pdfauthor={}}

\begin{document}
\maketitle
\begin{CornellOverviewBox}{Overview}
\textbf{Classification:} Standard library - Normative\\[2pt]
\textbf{Learning objective:} Understand the rules, terminology, and semantic consequences associated with Operations on atomic types.
\end{CornellOverviewBox}\begin{CornellSummaryBox}{Summary}
{\sffamily\bfseries\color{Accent} Summary}\par\vspace{3pt}Section 33.5.8.2 focuses on Operations on atomic types. Apply the specified interface together with its mandates, effects, result conditions, exception behavior, and complexity constraints. When applying it, preserve the distinction between mandatory requirements, recommendations, permissions, and behavior deliberately left undefined, unspecified, or implementation-defined.
\end{CornellSummaryBox}

\end{document}
